\documentclass[12pt]{article}
\usepackage[utf8]{inputenc}
\usepackage[T1]{fontenc}
\usepackage{hyperref}
\title{Bukiyip Language Resources for K-12 Curriculum Development}
\author{Compiled by Claude Sonnet 5.5 (claude-sonnet-5-5) for Andr\'as Kornai}
\date{}
\begin{document}
\maketitle

\section{Introduction}
Bukiyip (ISO 639-3: ape), also called Mountain Arapesh, is a Torricelli language of East Sepik Province, Papua New Guinea. It is distinct from Abelam, a Ndu language with which it is sometimes confused in search results. Only one substantial resource was found.

\section{Primary Educational Resources}
No primers, textbooks or school programmes were found in search results.

\section{Linguistic and Reference Materials}
\begin{itemize}
\item An Outline of Bukiyip Grammar (Robert J. Conrad and Kepas Wogiga, Pacific Linguistics C-113, 1991), openly available from the Australian National University repository: \url{https://openresearch-repository.anu.edu.au/items/d4478933-b15f-4283-a54f-8ceb6ce28c89}
\item Overview of the language (encyclopedic, secondary source): \url{https://en.wikipedia.org/wiki/Bukiyip_language}
\end{itemize}
No Bukiyip dictionary was found in search results.

\section{Religious and Community Resources}
None found in search results.

\section{Digital and Audio Resources}
No audio, video or corpus resources were found in search results.

\section{Curriculum Development Context}
The available documentation is a single descriptive grammar. A curriculum would first need a vocabulary and orthography documentation effort and community involvement; nothing beyond the grammar could be verified.

\end{document}